%%
%% winter-lecture26.tex
%% 
%% Made by Alex Nelson <pqnelson@gmail.com>
%% Login   <alex@lisp>
%% 
%% Started on  2026-03-05T08:09:43-0800
%% Last update 2026-03-05T08:09:43-0800
%% 

\lecture{}

% lecture 25
\begin{lemma}\label{lemma:19.4}
In the AKLB setup, for any $0\neq Q\in\Spec(B)$ and $P=Q\cap A$, we
have
\begin{equation*}
\dim_{A/P}(B/Q^{e_{Q}})=e_{Q}f_{Q}.
\end{equation*}
\end{lemma}

% The proof was in lecture 26
\begin{proof}
First, for any $i,j\geq0$, we show that $Q^{i}/Q^{i+1}\iso Q^{j}/Q^{j+1}$
as $(A/P)$-vector spaces.

\textsc{Step 1: Construct the morphism and prove it is injective}. For any $\pi\in Q\setminus Q^{2}$, the $(A/P)$ linear map
\begin{equation}
\begin{split}
\varphi\colon&B/Q\to Q^{i}/Q^{i+1}\\
&x+Q\mapsto x\pi^{i}+Q^{i+1}
\end{split}
\end{equation}
is clearly injective (since the preimage of zero is trivial).

\textsc{Step 2: Surjectivity}.
We will show $\varphi$ is surjective: consider the ideal (of $B$)
\begin{equation}
\pi B=\prod_{k}Q_{k}^{e_{k}}\subset Q,
\end{equation}
there exists a $k$ such that $Q_{k}=Q$ and $e_{k}=1$. Then
\begin{equation}
\Im(\varphi)=\pi^{i}B+Q^{i+1}
=Q^{i}(Q+\prod_{Q_{k}\neq Q}Q_{k}^{e_{k}})
\end{equation}
and we have one of two cases:

\textsc{Case 1:} $\prod_{Q_{k}\neq Q}Q_{k}^{e_{k}}=0$. Then $\pi B=Q$.
Hence $\varphi$ is clearly surjective.

\textsc{Case 2:} $\prod_{Q_{k}\neq Q}Q_{k}^{e_{k}}\neq0$.
Then
\begin{equation}
Q+\prod_{Q_{k}\neq Q}Q_{k}^{e_{k}}=B,
\end{equation}
and clearly $\varphi$ is surjective.

\textsc{Step 3: Check dimensions}.
We see
\begin{subequations}
  \begin{align}
\dim_{A/P}(Q^{i}/Q^{i+1}) &= \dim(B/Q)\\
&= f_{Q},
  \end{align}
\end{subequations}
and also
\begin{subequations}
  \begin{align}
\dim_{A/P}(B/Q^{e_{Q}})
&=\sum^{e_{Q}-1}_{i=0}\bigl(\dim_{A/P}(Q^{i}/Q^{e_{Q}})-\dim_{A/P}(Q^{i+1}/Q^{e_{Q}})\bigr)\\
&=\sum^{e_{Q}-1}_{i=0}\dim_{A/P}(Q^{i}/Q^{i+1})\\
&=e_{Q}f_{Q},
  \end{align}
\end{subequations}
as desired.
\end{proof}

\begin{corollary}\label{cor:19.5}
In the same setup as Lemma~\ref{lemma:19.4}, we have
\begin{equation*}
\sum_{Q\divides PB}e_{Q}f_{Q}=[L:K].
\end{equation*}
\end{corollary}

\begin{proof}
By Chinese remainder theorem, since
\begin{equation}
PB=\prod_{Q\divides PB}Q^{e_{Q}},
\end{equation}
we have
\begin{equation}
B/PB=\bigoplus_{Q\divides PB}B/Q^{e_{Q}}.
\end{equation}
By Lemmas~\ref{lemma:19.3} and \ref{lemma:19.4}, the claim holds.
\end{proof}

\begin{fact}\label{fact:math120b:winter2026:lecture26:fact1}
For ideals $I$ and $J$ of a ring $A$, we see that
\begin{equation*}
\frac{A/I}{J(A/I)}\iso\frac{A/J}{I(A/J)}\iso\frac{A}{I+J}.
\end{equation*}
\end{fact}
\begin{proof}[Proof hint]
By considering the canonical quotient morphisms $A/I\onto A/(I+J)$
and $A/J\onto A/(I+J)$.
\end{proof}

\begin{fact}\label{fact:math120b:winter2026:lecture26:fact2}
For any ideal $I\ideal A$, the canonical quotient morphism $A\onto A/I$
gives a bijection
\begin{equation*}
\left\{\begin{matrix}\mbox{ideals of }A\\
\mbox{containing }I\end{matrix}\right\}\stackrel{1:1}{\longleftrightarrow}\{\mbox{ideals of }A/I\}
\end{equation*}
This bijective correspondence preserves inclusion, prime ideals,
maximal ideals, etc.
\end{fact}

\begin{definition}
Let $A\subset B$ be a ring extension. The \define{Conductor} of $A$ in
$B$ is the ideal $C$ of $A$ (and of $B$) defined as
\begin{subequations}
\begin{equation}
C:=\{x\in A \mid xB\subset A\}
\end{equation}
which also satisfies
\begin{equation}
C=\{x\in B\mid xB\subset A\}.
\end{equation}
\end{subequations}
\end{definition}

\begin{lemma}\label{lemma:19.6}
In the AKLB setup, let $\theta\in B$ be such that $L=K(\theta)$ (by Theorem~\ref{thm:4.5}),
let $f$ be the minimal polynomial of $\theta$ over $A$ (so $f\in A[x]$),
let $C$ be the conductor of $A[\theta]$ in $B$. 
If prime ideal $P$ of $A$ is coprime to $C$ (so $P+C=A$),
then
\begin{enumerate}
\item $B/PB\iso A[\theta]/PA[\theta]$; and also
\item $B/PB\iso (A/P)[x]/(f(A/P)[x])$.
\end{enumerate}
\end{lemma}

\begin{proof}
\begin{enumerate}
\item Let $i\colon A[\theta]\into B$ be the inclusion morphism and
  $q_{PB}\colon B\onto B/PB$ be the canonical quotient map. We define
  $\phi=q_{PB}\circ i\colon A[\theta]\to B/PB$.

  We claim $\phi$ is surjective. This is because
\begin{subequations}
  \begin{align}
B&=PB + CB\\
\intertext{and}
CB&=C\subset A\subset A[\theta]\\
\intertext{gives us}
B&=P\theta+A[\theta].
  \end{align}
\end{subequations}
Hence $\phi$ is surjective.

We claim $\phi$ is injective.
Since
\begin{equation}\label{eq:math120b:lec26:C-ker-contained-in-PC}
C\ker(\phi)\subset CBP=PC
\end{equation}
and
\begin{equation}\label{eq:math120b:lec26:PC-contained-in-PA-theta}
PC\subset PA[\theta],
\end{equation}
we see
\begin{subequations}
  \begin{align}
\ker(\phi)
&=(PA[\theta]+C)\ker(\phi)\\
&=PA[\theta]+C\quad\mbox{by Eq~\eqref{eq:math120b:lec26:C-ker-contained-in-PC}}\\
&=PA[\theta]\quad\mbox{by Eq~\eqref{eq:math120b:lec26:PC-contained-in-PA-theta}}
  \end{align}
\end{subequations}
Hence
\begin{equation}
B/PB\iso A[\theta]/\ker(\phi)\iso A[\theta]/PA[\theta].
\end{equation}
\item The result follows from applying Fact~\ref{fact:math120b:winter2026:lecture26:fact1}
to the case where $I=(f(x))$ is principal, $J=PA[x]$, and the fact
$A[\theta]\iso A[x]/(f(x))$. \qedhere
\end{enumerate}
\end{proof}

\begin{theorem}[Dedekind--Kummer]\label{thm:19.7}
In the setup of Lemma~\ref{lemma:19.6}, let 
\begin{subequations}
\begin{equation}
f(x)\equiv \varphi_{1}(x)^{e_{1}}(\cdots)\varphi_{r}(x)^{e_{r}}\bmod{P}
\end{equation}
(i.e., the prime factorization in $(A/P)[x]$), each
$\varphi_{i}\in(A/P)[x]$ is rreducible. Then
\begin{equation}
Q_{i}=\varphi_{i}(\theta)B+PB
\end{equation}
are prime ideals of $B$ and
\begin{equation}
PB=Q_{1}^{e_{1}}(\cdots)Q_{r}^{e_{r}}
\end{equation}
is the prime factorization of $PB$ in $B$. Moreover,
\begin{equation}
f_{Q_{i}}=\deg(\varphi_{i})
\end{equation}
for each $i$.
\end{subequations}
\end{theorem}

\begin{proof}
For
\begin{equation}
\psi\colon B\onto B/PB
\end{equation}
and, by the last lemma, the codomain is isomorphic to
\begin{equation}
\psi\colon B\onto\frac{(A/P)[x]}{(f(A/P))[x]},
\end{equation}
the prime ideal of the codomain generated by $(\overline{\varphi_{i}(x)})$
corresponds to
\begin{equation}
Q_{i}=\varphi_{i}(A)B+PB
\end{equation}
by Fact~\ref{fact:math120b:winter2026:lecture26:fact2}. Since
\begin{equation}
Q_{1}^{e_{1}}(\cdots)Q_{r}^{e_{r}}\subset\ker(\psi),
\end{equation}
we have
\begin{equation}
Q_{1}^{e_{1}}(\cdots)Q_{r}^{e_{r}}\subset PB
\end{equation}
and thus
\begin{equation}
PB=Q_{1}^{e'_{1}}(\cdots)Q_{r}^{e'_{r}}
\end{equation}
where each $e'_{i}\leq e_{i}$. Then
\begin{subequations}
  \begin{align}
f_{i} &= \dim_{A/P}(B/Q_{i})\\
&= \dim_{A/P}\left(\frac{(A/P)[x]}{\varphi_{i}(x)((A/P)[x])}\right)\\
&= \deg(\varphi_{i}),
  \end{align}
\end{subequations}
and so by Corollary~\ref{cor:19.5},
\begin{subequations}
  \begin{align}
[L:K]
&=\sum^{r}_{i=1}e'_{i}f_{i}\\
&\leq\sum^{r}_{i=1}e_{i}f_{i}=\deg(f)=[L:K],
  \end{align}
\end{subequations}
which forces $e'_{i}=e_{i}$ for each $i$.
\end{proof}

\begin{definition}
In the AKLB setup, let $\hom_{K}(L,\closure{K})=\{\sigma_{1},\dots,\sigma_{n}\}$
and let $\{b_{1},\dots,b_{n}\}\subset B$ be a basis of $L/K$ (called
an \define{Integral Basis}). We define the \define{Discriminant}
$D_{L/K}$ to be the ideal generated by
\begin{equation}
\{d(b_{1},\dots,b_{n})=\det\bigl((\sigma_{i}(b_{j}))\bigr)^{2}\mid
(b_{1},\dots,b_{n})\mbox{ is an integral basis}\}.
\end{equation}
\end{definition}

\begin{theorem}[Dedekind's discriminant theorem]\label{thm:19.8}
In the AKLB setup, a prime ideal $P$ of $A$ is unramified (\S\ref{lemma:19.2})
if and only if $P\ndivides D_{L/K}$. In particular, there exists only
finitely many ramified prime ideals.
\end{theorem}

\begin{proof}
\forwardproof\ Homework.

\backwardproof\ Assume $P\ndivides D_{L/K}$. By the proof of
Proposition~\ref{prop:17.3}, for $\alpha\in B$ such that $L=K(\alpha)$
and
\begin{subequations}
  \begin{align}
D & := d(1,\alpha,\dots,\alpha^{n-1})\\
&=\prod_{1\leq i<j\leq n}(\sigma_{j}(\alpha)-\sigma_{i}(\alpha))^{2}\\
&=\det\begin{pmatrix}
\sigma_{1}(1) & \sigma_{1}(\alpha) & \dots & \sigma_{1}(\alpha)^{n-1}\\
\sigma_{2}(1) & \sigma_{2}(\alpha) & \dots & \sigma_{2}(\alpha)^{n-1}\\
\vdots & \vdots & \ddots & \vdots\\
\sigma_{n}(1) & \sigma_{n}(\alpha) & \dots & \sigma_{n}(\alpha)^{n-1}
\end{pmatrix}^{2}.
  \end{align}
\end{subequations}
Then
\begin{equation}
DB\subset A[\theta],
\end{equation}
that is to say, the conductor $C$ of $A[\theta]$ in $B$ divides
$D$. In particular, $P$ and $C$ are corpime since $C\divides D$ but
$P\ndivides D$. Then by Theorem~\ref{thm:19.7}, we have
\begin{equation}
PB=\prod_{Q\divides PB}Q.
\end{equation}
Also, by the proof of Lemma~\ref{lemma:19.6}, $A[\theta]\onto B/PB$ is
surjective, and so
\begin{equation}
B/PB\iso\bigoplus_{Q\divides PB}B/Q
\end{equation}
is generated by $\theta\bmod{P}$, i.e., it is generated by the set
\begin{equation}
\{\alpha\bmod{Q}\mid Q\mbox{ divides }PB\}.
\end{equation}
Since $f$ (and thus $\alpha$) is separable, then so is $(B/Q)/(A/P)$.
In other words, $P$ is unramified.
\end{proof}